\section{Conclusion}
\label{sec:conclusion}
AdaSolver turns spatial sampling density from a fixed discretization choice into a controllable design variable at test time. By compensating for the induced measure shift, it decouples where computation is allocated from the reference operator being approximated, establishing spatial sampling as a new axis of test-time scaling for point-based neural operators. Our experiments further show that effective sampling strategies can arise directly from computational-physics knowledge: classical refinement heuristics, geometric and flow features, and solution-dependent structures can all be translated into improved accuracy–compute trade-offs, and ultimately into faster surrogate-based aerodynamic design. More broadly, this perspective suggests a path toward measure-aware adaptive computation across a wider class of transformer-based neural operators, as well as training-time augmentation over diverse sampling distributions, allowing neural operators to scale not only in model capacity or iterative computation, but also in where computation is spent in physical space.